\scanpage{111}

donc l'autom. de $\hat\Sigma_{\mathcal M}$ induit par $v\mapsto v'$ est aussi l'antipodisme, induit par la réflexion $\underset{\substack{\Vert\\v}}{u+t.1}\longmapsto\underset{\substack{\Vert\\-v'}}{u-t.1}$. Ceci

\deleted{\unsure{par transport}}

achève de décrire notre iso, $(*)$, \unsure{satisfais.|satisfaisant} à a) c) d) il reste à vérifier b), ou ce qui revient au $=$, la formule

\[
f_{-\omega}(x,y)=\sigma f_\omega(x,y),
\]

Ce qui résulte du fait que $(\omega,x,y)\mapsto f_\omega(x,y)$, $X\times X\times\Omega\longrightarrow\hat\Sigma$, ayant été défini \unsure{fonctoriellt|fonctoriellement} en $(E,q,\omega)$, commute à l'action de l'isom

\[
-\operatorname{id}_E:(E,q,\omega)\mathrel{\unsure{\longrightarrow|\text{---}}}(E,q,-\omega)
\]

\figure{\begin{tikzpicture}[x=1cm,y=1cm]\draw (0,0) .. controls (.26,.06) and (.41,.03) .. (.59,.08) .. controls (.85,.02) and (1.05,.17) .. (1.25,.12) .. controls (1.47,.08) and (1.68,.13) .. (1.9,.16);\end{tikzpicture}}

Il reste à faire la jonction avec les développements du \S~1, relatifs au cas où on se donne, en $=$ temps que $q$ \inserted{sur $V$} définissant la quadrique affine d'équation $q(x)+1=0$, une $\omega_0\in\Gamma\det E$ satisfaisant (non pas $\delta'(q)=\omega_0^{\otimes-2}$ mais\ldots)

\[
\delta'(q)=-\omega_0^{\otimes-2}\qquad\text{i.e.}\quad\delta'(-q)=\omega_0^{\otimes(-2)}
\]

% archive sheet number: 110
